% Glossary entries (must be defined in the preamble for no-index mode).
% First body mention of each term uses \gls{<label>}, which renders the term
% and hyperlinks back to the Glossary.
\newglossaryentry{dtu}{
  name={Differential Transcript Usage (DTU)},
  first={Differential Transcript Usage (DTU)},
  text={DTU},
  description={Differential transcript usage occurs when the relative expression of multiple transcripts arising from the same gene changes between different conditions \cite{young2024dtu}. It is independent of changes in the gene's total expression, so it is not captured by gene-level differential expression analysis and requires transcript-level resolution to detect.}
}

\newglossaryentry{deg}{
  name={differential gene expression (DGE)},
  first={differential gene expression (DGE)},
  text={DGE},
  description={The analysis of differences in a gene's total expression---the summed expression of all its isoforms---between groups of samples. It measures overall gene output and is blind to changes in the relative usage of individual isoforms, which are instead captured by differential transcript usage.}
}

\newglossaryentry{threePrimeEnd}{
  name={3' end},
  description={The terminal end of an RNA molecule at which transcription finishes; in mRNA, it is where the poly-A tail is added.}
}

\newglossaryentry{threePrime}{
  name={3'-biased sequencing},
  description={A property of some scRNA-seq protocols (such as 10x Genomics Chromium) in which sequencing reads are derived only from the 3' terminus of each transcript. Because reading quality decreases by distance and because capture is anchored at the transcript's 3' end, coverage does not extend into the transcript body, so internal splice junctions are never observed. Only isoforms that differ at or near their 3' ends are therefore distinguishable \cite{regan2022practical}.}
}

\newglossaryentry{isoformSwitch}{
  name={isoform switching},
  description={A change between conditions in the relative usage of a gene's isoforms, such that the fraction of the gene's expression contributed by one isoform increases while that of another decreases. It reflects altered alternative splicing or transcription and is quantified by the change in each isoform's fraction of total gene expression \cite{vitting2017isoformswitch}.}
}

\newglossaryentry{sparsity}{
  name={sparsity},
  description={The very high proportion of zero counts in single-cell RNA-seq data. Zero fractions reach approximately 90\,\% of matrix entries in droplet-based data, compared with roughly 10--40\,\% in bulk RNA-seq. These zeros arise from both the true absence of a transcript (biological zeros) and technical failures (dropouts), which are not distinguishable in the raw counts \cite{jiang2022statistics}.}
}

\newglossaryentry{dropout}{
  name={technical dropout},
  plural={technical dropouts},
  description={A zero count produced by a technical limitation of the assay---failed capture, reverse transcription, or sequencing---rather than by the true absence of the transcript. It is also described as a non-biological zero. Because dropouts are numerically identical to biological zeros, they cannot be identified directly in the raw count matrix \cite{jiang2022statistics}.}
}

\newglossaryentry{zeroInflated}{
  name={zero-inflated distribution},
  plural={zero-inflated distributions},
  description={A statistical model for count data that is a mixture of a point mass at zero and a standard count distribution (such as Poisson or negative binomial). It is used when the data contain more zeros than the count distribution alone predicts, treating those excess zeros as generated by a separate process. In single-cell RNA-seq the term has been applied inconsistently, and its use is debated because a zero may reflect either measurement error or genuine absence of expression \cite{sarkar2021separating}.}
}

% No extra period between a description and its page list.
\renewcommand*{\glspostdescription}{}

\newglossaryentry{isoform}{
  name={isoform},
  plural={isoforms},
  description={A distinct mature transcript variant produced from the same gene by alternative splicing.}
}

\newglossaryentry{scRnaSeq}{
  name={single-cell RNA sequencing (scRNA-seq)},
  first={Single-cell RNA sequencing (scRNA-seq)},
  text={scRNA-seq},
  description={RNA sequencing performed on thousands of individual cells separately, exposing per-cell expression.}
}

\newglossaryentry{bulkRnaSeq}{
  name={bulk RNA-seq},
  description={Sequencing of RNA pooled from many cells at once, which averages away cell-to-cell differences.}
}

\newglossaryentry{fullLength}{
  name={full-length sequencing},
  description={A low-throughput, higher-cost protocol that captures complete RNA molecules, enabling precise isoform resolution.}
}

\newglossaryentry{read}{
  name={read},
  plural={reads},
  description={A short nucleotide sequence produced by the sequencing machine, representing the raw data of RNA-seq.}
}

\newglossaryentry{rnaSeq}{
  name={RNA sequencing (RNA-seq)},
  description={A technique that determines the sequence and abundance of RNA molecules in a sample; it is the basis of the single-cell protocols discussed in this thesis.}
}

\newglossaryentry{splicing}{
  name={alternative splicing},
  description={The regulated joining of different exon combinations from one gene, producing multiple distinct isoforms.}
}

\newglossaryentry{mutualExclusivity}{
  name={mutual exclusivity},
  description={Two features that rarely appear together in the same cell (a negative association), as opposed to co-expression. For isoforms of one gene, mutual exclusivity is consistent with alternative splicing or isoform switching across cell states, and it is the signal COTAN screens for.}
}

\newglossaryentry{gdi}{
  name={Global Differentiation Index (GDI)},
  first={Global Differentiation Index (GDI)},
  text={GDI},
  description={A cluster-specificity measure used to identify marker genes and validate clustering quality.}
}

\newglossaryentry{coexpression}{
  name={co-expression},
  description={The tendency of two features (genes or transcripts) to be detected together across cells. In sparse single-cell data, raw co-expression measured by correlation can reflect shared zeros rather than shared regulation, which is the problem COTAN addresses by testing presence/absence independence.}
}

\newglossaryentry{contingencyTable}{
  name={contingency table},
  description={A $2 \times 2$ table counting, across all cells, the four presence/absence combinations of two features; it tests whether they co-occur more or less than chance.}
}

\newglossaryentry{pseudoBulk}{
  name={pseudo-bulk aggregation},
  description={The practice of summing read counts across groups of single cells to simulate bulk RNA-seq depth, which recovers statistical power at the expense of single-cell heterogeneity.}
}

\newglossaryentry{orthogonalValidation}{
  name={orthogonal validation},
  description={Independent confirmation of a computational result by a different experimental technique.}
}
% Suppress the package-generated heading; the chapter provides its own.
\renewcommand{\glossarysection}[2][]{}
% Print cue: mark glossary terms at first use so paper readers spot them.
\renewcommand{\glstextformat}[1]{\textit{#1}}
